\section*{Hoofdstuk \ref{ch:b1:first-law} --- De eerste hoofdwet van de thermodynamica}

\begin{solution}{exo:b1:first-law:1}
$Q = C_{P,m}\Delta T = 29.1 \times 100 = \qty{2.91}{kJ}$; $W = -nR\Delta T =
\qty{-0.83}{kJ}$; $\Delta U = Q + W = \qty{2.08}{kJ} = C_{V,m}\Delta T$.
\end{solution}

\begin{solution}{exo:b1:first-law:2}
$W = -nRT\ln(V_2/V_1) = -2494\ln(0.25) = \qty{+3.46}{kJ}$; $\Delta U = 0$, dus
$Q = \qty{-3.46}{kJ}$ (afgestaan aan het bad).
\end{solution}

\begin{solution}{exo:b1:first-law:3}
$T_2 = T_1 \times 10^{\gamma - 1} = 300 \times 10^{0.4} = \qty{754}{K}$; $P_2 =
10^{1.4} = \qty{25}{bar}$.
\end{solution}

\begin{solution}{exo:b1:first-law:4}
Argon: $C_{V,m} = \tfrac32 R = 12.5$, $C_{P,m} = \tfrac52 R = 20.8$ J/(\omterm{def:b1:kinetic-theory:scales}{mol} K),
$\gamma = 1.67$; stikstof: $20.8$, $29.1$, $\gamma = 1.40$. Voor water is de
uitzetting bij verwarming minuscuul, zodat de arbeid $P\,\dd V$ bij constante druk
verwaarloosbaar is tegen de \omterm{def:b1:first-law:heat}{warmte}: $c_P - c_V \approx 0$.
\end{solution}

\begin{solution}{exo:b1:first-law:5}
$T_f = (200 \times 80 + 300 \times 20)/500 = \qty{44}{\degreeCelsius}$. Met de
beker: $(200 \times 4.18 \times 80 + (300 \times 4.18 + 80) \times 20)/(500 \times
4.18 + 80) = \qty{43}{\degreeCelsius}$.
\end{solution}

\begin{solution}{exo:b1:first-law:6}
IJs: opwarmen tot $0$: $0.050 \times 2100 \times 10 = \qty{1.05}{kJ}$; smelten:
\qty{16.7}{kJ}; samen \qty{17.8}{kJ}. Het water kan \qty{25.1}{kJ} afstaan tot
$0$: alles smelt; dan $0.250 \times 4180\,T_f = 25.1 - 17.8$ kJ: $T_f =
\qty{7}{\degreeCelsius}$. Met \qty{150}{g}: nodig \qty{53.3}{kJ} $>$ \qty{25.1}{kJ}:
er smelt maar $(25.1 - 3.15)/334 = \qty{66}{g}$ ijs, $T_f = \qty{0}{\degreeCelsius}$,
\qty{84}{g} ijs blijft over.
\end{solution}

\begin{solution}{exo:b1:first-law:7}
$W = -P_1(V_2 - V_1) + 0 - P_2(V_1 - V_2) + 0 = (P_2 - P_1)(V_2 - V_1) > 0$:
het kringproces wordt tegen de klok in doorlopen (uitzetten bij lage druk,
samendrukken bij hoge), dus \emph{ontvangt} het gas netto arbeid, gelijk aan de
ingesloten oppervlakte; met de klok mee zou het haar leveren --- een motor.
\end{solution}

\begin{solution}{exo:b1:first-law:8}
$T_2 = 300 \times 7^{0.286} = \qty{523}{K}$. De hete samengeperste lucht verwarmt de
cilinder bij elke slag door geleiding. In de band koelt de lucht bij vast
volume tot de omgevingstemperatuur af: $P \propto T$, de druk zakt met de verhouding
van de temperaturen --- even bijpompen.
\end{solution}

\begin{solution}{exo:b1:first-law:9}
$W = 0$ (geen uitwendige druk om tegenin te duwen), $Q = 0$: $\Delta U = 0$,
en dus $\Delta T = 0$ voor een \omterm{def:b1:kinetic-theory:model}{ideaal gas}. Onomkeerbaar: het gas stroomt nooit
vanzelf terug. Een echt gas, waarvan de moleculen elkaar aantrekken, koelt licht af
(een deel van de \omterm{thm:b1:work-and-energy:ke}{kinetische energie} wordt \omterm{def:b1:work-and-energy:conservative}{potentiële energie} als zij uit elkaar gaan).
\end{solution}

\begin{solution}{exo:b1:first-law:10}
$Q = L_v = \qty{2.26}{MJ}$; $W = -P(V_g - V_l) = -1.013\times10^5 \times 1.67 =
\qty{-0.17}{MJ}$; $\Delta U = \qty{2.09}{MJ}$: $93\%$ van de \omterm{def:b1:first-law:heat}{warmte} gaat naar het
uit elkaar halen van de moleculen (\omterm{def:b1:kinetic-theory:U}{inwendige energie}), $7\%$ naar het wegduwen van de
atmosfeer.
\end{solution}

\begin{solution}{exo:b1:first-law:11}
$W = -P_2(V_2 - V_1)$, $\Delta U = nC_{V,m}(T_2 - T_1)$, $V_2 = nRT_2/P_2$, $V_1
= nRT_1/P_1$: $\tfrac52(T_2 - 300) = -T_2 + 300 \times 5$, $3.5T_2 = 2250$,
$T_2 = \qty{643}{K}$, $V_2 = \qty{10.7}{L}$. Omkeerbaar: $T_2 = 300 \times
5^{0.286} = \qty{475}{K}$, $V_2 = \qty{7.9}{L}$. De plotselinge compressie verricht
meer arbeid op het gas (de volle \qty{5}{bar} werkt vanaf het begin) en eindigt
warmer en minder samengedrukt.
\end{solution}

\begin{solution}{exo:b1:first-law:12}
$\rho = PM/RT = \qty{1.20}{kg/m^3}$: $\sqrt{\gamma P/\rho} = \sqrt{1.4 \times
1.013\times10^5/1.20} = \qty{343}{m/s}$; $\sqrt{P/\rho} = \qty{290}{m/s}$, $15\%$
te laag. Met $P/\rho = RT/M$: $c = \sqrt{\gamma RT/M}$; tegen $u = \sqrt{3RT/M}$
= \qty{502}{m/s}: $c/u = \sqrt{\gamma/3} = 0.68$.
\end{solution}

\begin{solution}{pb:b1:first-law:1}
\textbf{1.} $n = P_1V_1/RT_1 = 10^5 \times 2.12\times10^{-4}/2494 = \qty{8.5e-3}{mol}$.

\textbf{2.} $PV^\gamma$ constant: $V_2 = V_1(P_1/P_2)^{1/\gamma} = 212 \times
7^{-0.714} = \qty{53}{cm^3}$.

\textbf{3.} $T_2 = T_1(P_2/P_1)^{(\gamma - 1)/\gamma} = 300 \times 7^{0.286} =
\qty{523}{K}$.

\textbf{4.} $W = nC_{V,m}(T_2 - T_1) = 8.5\times10^{-3} \times 20.8 \times 223 =
\qty{39}{J}$.

\textbf{5.} Toe te voegen lucht: $\Delta n = \Delta P\,V/RT = 6\times10^5 \times 2\times10^{-3}/
2494 = \qty{0.48}{mol}$: $0.48/8.5\times10^{-3} = 57$ slagen, ongeveer
\qty{2.2}{kJ} --- een minuut eerlijk werk.

\textbf{6.} Bij constant volume geldt $P \propto T$: de druk zakt terwijl de lucht
afkoelt (tot $300/523$ voor de lucht van de laatste slag); de fietser geeft na
een pauze nog een paar slagen.

\textbf{7.} $W_{\text{iso}} = nRT\ln(P_2/P_1) = 8.5\times10^{-3} \times 2494 \times
1.95 = \qty{41}{J} > \qty{39}{J}$: de isotherme weg drukt het gas verder
samen (tot \qty{30}{cm^3} in plaats van \qty{53}{cm^3}) om \qty{7}{bar} te bereiken,
omdat het koel blijft --- meer volume doorlopen, meer arbeid.

\textbf{8.} $n = 10^5 \times 5\times10^{-4}/2494 = \qty{0.020}{mol}$; $T_2 =
300 \times 20^{0.4} = \qty{994}{K}$.

\textbf{9.} $P_2 = 20^{1.4} = \qty{66}{bar}$.

\textbf{10.} \qty{994}{K} ligt ver boven het zelfontbrandingspunt van de brandstof: de
ingespoten brandstof ontsteekt bij aanraking. Een benzinemotor perst zijn
lucht--brandstofmengsel maar tot ongeveer \qty{750}{K} samen, zodat het niet
vóór de vonk ontbrandt; een te hoge verhouding geeft ``pingelen''.

\textbf{11.} $W = nC_{V,m}(T_2 - T_1) = 0.020 \times 20.8 \times 694 = \qty{289}{J}$.

\textbf{12.} $25 \times 289 = \qty{7.2}{kW}$ per cilinder, opgeslagen in het hete
gas en teruggegeven bij de expansie.

\textbf{13.} De onomkeerbaarheid (het gas volgt de quasistatische weg niet;
er wordt extra arbeid gedissipeerd) verhoogt $T_2$; warmteverlies aan de wanden
verlaagt haar. In een echte motor wint het tweede meestal: $T_2$ ligt iets
onder de ideale waarde.

\textbf{14.} $Q = 20\times10^{-6} \times 43\times10^6 = \qty{860}{J}$.

\textbf{15.} Isobaar: $Q = nC_{P,m}(T_3 - T_2)$: $T_3 - T_2 = 860/(0.020 \times
29.1) = \qty{1480}{K}$, $T_3 = \qty{2470}{K}$; $V_3 = V_2T_3/T_2 = 25 \times
2.48 = \qty{62}{cm^3}$.

\textbf{16.} $W = -P_2(V_3 - V_2) = -66\times10^5 \times 37\times10^{-6} =
\qty{-244}{J}$: door het gas geleverde arbeid.

\textbf{17.} $T_4 = T_3(V_3/V_1)^{\gamma - 1} = 2470 \times (0.124)^{0.4} =
\qty{1070}{K}$; $W = nC_{V,m}(T_4 - T_3) = 0.020 \times 20.8 \times (-1400) =
\qty{-582}{J}$.

\textbf{18.} Netto geleverde arbeid $= 244 + 582 - 289 = \qty{537}{J}$; verhouding
$537/860 = 62\%$ (het ideale dieselrendement bij deze instellingen; echte
motoren halen ongeveer $40\%$).

\textbf{19.} $1 - 300/2470 = 88\%$: het ideale dieselkringproces blijft ruim
onder de grens van Carnot, omdat het zijn \omterm{def:b1:first-law:heat}{warmte} over een reeks
temperaturen ontvangt en niet bij de hoogste.

\textbf{20.} Over een kringproces is $\Delta U = 0$: $Q_{\text{uit}} = -(860 - 537) =
\qty{-323}{J}$, dus \qty{323}{J} verdwijnt met de uitlaat; controle:
$nC_{V,m}(T_4 - T_1) = 0.020 \times 20.8 \times 770 = \qty{320}{J}$.

\textbf{21.} De trilling is snel (een seconde) vergeleken met warmtegeleiding
door \qty{10}{L} lucht: adiabatisch. Linearisering van
$PV^\gamma$ geeft $\dd P/P = -\gamma\,\dd V/V = -\gamma Sx/V$.

\textbf{22.} Kracht op de kogel $S\,\dd P = -\gamma PS^2x/V$: $m\ddot x =
-(\gamma PS^2/V)x$, harmonisch, $\omega_0^2 = \gamma PS^2/mV$.

\textbf{23.} $\omega_0^2 = 1.4 \times 10^5 \times 4\times10^{-8}/(0.0165 \times 0.010)
= \qty{33.9}{s^{-2}}$: $T = 2\pi/5.83 = \qty{1.08}{s}$. Uit $T = \qty{1.10}{s}$:
$\gamma = 4\pi^2mV/(T^2PS^2) = 4\pi^2 \times 0.0165 \times 0.010/(1.21 \times 10^5
\times 4\times10^{-8}) = 1.35$.

\textbf{24.} Wrijving dempt de beweging en verschuift, als zij droog is, het
evenwicht; lekkage langs de kogel laat lucht ontsnappen tijdens een compressie,
wat de terugdrijvende kracht verkleint: beide verlengen de periode en drukken
$\gamma$ omlaag --- zoals de $1.35$ hierboven.

\textbf{25.} Pomp en diesel: $Q = 0$, dus $\Delta U = W$ --- arbeid werd
\omterm{def:b1:kinetic-theory:U}{inwendige energie} en temperatuur, volgens de \omterm{prop:b1:first-law:transformations}{wet van Laplace} met exponent
$\gamma$. Kogel: een snelle compressie is adiabatisch, en $\gamma$ bepaalt de
stijfheid van een luchtveer.
\end{solution}
